\TieudeT{PHONG TỎA VẬT LÍ}
\subsubsection*{PHẦN I. Thí sinh trả lời câu hỏi từ câu 1 đến câu 18. Mỗi câu hỏi thí sinh chỉ chọn một phương án}
\setcounter{ex}{0}
%mac
%\loai{Lực}

\begin{ex}%book%[1P2Y1-1][Loai1]
Một lực được biểu diễn bằng
\choice
{một đường thẳng}
{một đoạn thẳng}
{một mũi tên}
{\True một vectơ}
%	\loigiai{
%		Một lực được biểu diễn bằng một vectơ, gồm có phương và chiều.
%	}
\haidong{7} \end{ex}


%\begin{ex}%book
%		Phát biểu nào sau đây đúng?
%	\choice
%	{\True Khi một vật bị biến dạng hoặc vận tốc của vật thay đổi thì chắc chắn đã có lực tác dụng lên vật}
%	{Khi một vật đang chuyển động mà đột nhiên không còn lực nào tác dụng lên vật nữa thì vật sẽ dừng lại ngay lập tức}
%	{Lực là nguyên nhân gây ra chuyển động vì khi ta tác dụng lực lên một vật đang đứng yên thì khi đó vật bắt đầu chuyển động}	
%{Theo định luật I Newton, nếu một vật không chịu tác dụng của lực nào thì vật phải đứng yên}
%	\loigiai{
%}
% \haidong{7} \end{ex}
\begin{ex}%book%[1P2Y3-1][Loai1]
Lực \textbf{không} phải là nguyên nhân làm cho
\choice
{hướng chuyển động của vật thay đổi}
{hình dạng của vật thay đổi}
{độ lớn vận tốc của vật thay đổi}
{\True vật chuyển động}
\loigiai{
Lực không gây ra chuyển động cho vật, nó chỉ làm vật biến dạng hoặc thay đổi chuyển động của vật.
}
\haidong{7} %<MyLT1>
\end{ex}
\begin{ex}%book%[1P2B3-1][Loai1]
Phát biểu nào sau đây là \textbf{sai}? Lực là nguyên nhân làm cho
\choice
{\True vật chuyển động}
{hình dạng của vật thay đổi}
{độ lớn vận tốc của vật thay đổi}
{hướng chuyển động của vật thay đổi}
\loigiai{
Lực là nguyên nhân làm biến đổi chuyển động.
}
\haidong{7} \end{ex}
\begin{ex}%book%[1P2B3-1][Loai1]
Khi vật chỉ chịu tác dụng của một lực duy nhất thì nó
\choice
{chỉ biến dạng mà không biến đổi vận tốc}
{chuyển động thẳng đều mãi}
{chỉ biến đổi vận tốc mà không bị biến dạng}
{\True bị biến dạng hoặc biến đổi vận tốc}
\haidong{7} \end{ex}

%\loai{Giá của lực}
\begin{ex}%book
Giá của vectơ lực là
\choice
{đoạn thẳng mang vectơ lực}
{đoạn thẳng vuông góc vectơ lực}
{đường thẳng vuông góc vectơ lực}
{\True đường thẳng mang vectơ lực}
\haidong{7} \end{ex}
%\loai{Phân tích lực}
\begin{ex}%book
Phép phân tích lực cho phép ta thay thế
\choice
{một lực bằng một lực khác}
{\True một lực bằng hai hay nhiều lực thành phần}
{nhiều lực bằng một lực duy nhất}
{các vectơ lực bằng vectơ gia tốc}
\haidong{7} \end{ex}

%\newpage
%\loai{Quy tắc hình bình hành}
\begin{ex}%book
Khi có hai vectơ $\vec{F}_{1}$ và $\vec{F}_{2}$ đồng quy, tạo thành hai cạnh của một hình bình hành thì vectơ tổng hợp lực $\overrightarrow{F}$
\choice
{có điểm đặt tại một đỉnh bất kì của hình bình hành}
{\True có phương trùng với đường chéo của hình bình hành}
{có độ lớn $F=F_{1}+F_{2}$}
{cùng chiều với $\overrightarrow{F}_{1}$ hoặc $\overrightarrow{F}_{2}$}
\loigiai{Chọn B		
}\haidong{7} \end{ex}

\begin{ex}%book
Phát biểu nào sau đây về phép tổng hợp lực là \textbf{sai}?
\choice
{Xét về mặt toán học, tổng hợp lực là phép cộng các vectơ lực cùng tác dụng lên một vật}
{Lực tổng hợp có thể xác định bằng quy tắc hình bình hành}
{\True Độ lớn của lực tổng hợp luôn bằng tổng độ lớn của các lực thành phần tác dụng lên vật}
{Tổng hợp lực là phép thay thế các lực tác dụng đồng thời vào cùng một vật bằng một lực có tác dụng tương đương các lực thành phần}

\loigiai{Chọn B

}\haidong{7} \end{ex}

%\loai{Hợp lực}
\begin{ex}%book
Hai người cột hai sợi dây vào đầu một chiếc xe và kéo. Lực kéo xe lớn nhất khi hai lực kéo $\overrightarrow{F}_{1}$ và $\overrightarrow{F}_{2}$
\choice
{vuông góc với nhau}
{ngược chiều với nhau}
{\True cùng chiều với nhau}
{tạo với nhau một góc $45^{\circ}$}

\loigiai{Chọn C

}\haidong{7} %<MyLT1>
\end{ex}
\begin{ex}%book%[1P2Y1-1][Loai1]
Hai lực đồng quy $\vec{F}_1$, $\vec{F}_2$ hợp với nhau một góc $\alpha$, hợp lực của hai lực này có độ lớn bằng
\choice
{$F=\sqrt {F_1+F_2}$}
{\True $F=\sqrt {F_1^2+F_2^2+2F_1F_2\cos \alpha}$}
{$F=F_1^2+F_2^2$}
{$F=F_1-F_2$}
\loigiai{
Ta có: $F=\sqrt {F_1^2+F_2^2+2F_1F_2\cos \alpha}$.
}
\haidong{7} \end{ex}
\begin{ex}%book%[1P2Y1-1][Loai1]
Hai lực $\vec{F}_1$, $\vec{F}_2$ có độ lớn bằng nhau và hợp với nhau một góc $\alpha$. Hợp lực của chúng có độ lớn bằng
\choice
{$F=F_1-F_2$}
{$F=F_1+F_2$}
{\True $F=2F_1\cos\dfrac{\alpha}{2}$}
{$F=2F_1\cos\alpha$}
\loigiai{
Ta có: $F=2F_1\cos \dfrac{\alpha }{2}$.
}
\haidong{7} \end{ex}
%\loai{Điều kiện giới hạn của hợp lực}
\begin{ex}%book%[1P2Y1-1][Loai1]
Gọi $F_1$ và $F_2$ là độ lớn của hai lực thành phần, F là hợp lực của chúng. Trong mọi trường hợp
\choice
{F luôn luôn lớn hơn cả $F_1$ và $F_2$}
{F luôn luôn nhỏ hơn cả $F_1$ và $F_2$}
{\True F thỏa mãn $\left|{F_1-F_2}\right|\le F\le F_1+F_2$}
{F không bao giờ bằng $F_1$ hoặc $F_2$}
\loigiai{ 
Hợp lực phải thỏa mãn: $|F_1-F_2| \le F \le F_1+F_2$.
}	
\haidong{7} \end{ex}
\begin{ex}%book%[1P2Y1-1][Loai1]
Cho hai lực đồng quy có độ lớn $F_1=16\ N,\ F_2=12\ N$. Độ lớn hợp lực của chúng \textit{có thể} là
\choice
{\True $F=20\ N$}
{$F=30\ N$}
{$F=3,5\ N$}
{$F=2,5\ N$}
\loigiai{ 
Hợp lực phải thỏa mãn: $|F_1-F_2| \le F \le F_1+F_2$ hay $4\ N \le F \le 28\ N$.
}
\haidong{7} \end{ex}
\begin{ex}%book%[1P2Y1-1][Loai1]
Có hai lực đồng quy có cùng độ lớn $9\ N$ và $12\ N$. Trong số các giá trị sau đây, giá trị nào \textit{có thể} là độ lớn của hợp lực?
\choice
{$1\ N$}
{$2\ N$}
{\True $15\ N$}
{$25\ N$}
\loigiai{ 
Hợp lực phải thỏa mãn: $|F_1-F_2| \le F \le F_1+F_2$ hay $3\ N \le F \le 21\ N$.
}
\haidong{7} \end{ex}
\begin{ex}%book%[1P2Y1-1][Loai1]
Cho hai lực đồng quy có độ lớn bằng $5\ N$ và $10\ N$. Giá trị của hợp lực \textbf{không} thể là giá trị nào trong các giá trị sau đây?
\choice
{\True $16\ N$}
{$5\ N$}
{$12\ N$}
{$6\ N$}
%	\loigiai{
%		Ta có: $|F_1-F_2| \le F_{hl} \le F_1+F_2 \Rightarrow 5 \le F_{hl} \le 15 \Rightarrow$ Hợp lực không thể có thể có độ lớn bằng 16 N.
%	}
\haidong{7} \end{ex}
\begin{ex}%book
Cho hai lực đồng quy có độ lớn bằng $8\ N$ và $12\ N$. Giá trị của hợp lực \textbf{không} thể là giá trị nào trong các giá trị sau đây?
\choice
{$19\ N$}
{$4\ N$}
{$7\ N$}
{\True $21\ N$}

\loigiai{Chọn D

}
\haidong{7} \end{ex}
\begin{ex}%book
Hai lực có giá đồng quy có độ lớn $7\ N$ và $13\ N$. Độ lớn hợp lực của hai lực này \textbf{không} thể có giá trị nào sau đây?
\choice
{$7\ N$}
{$13\ N$}
{$20\ N$}
{\True $22\ N$}

\loigiai{Chọn D

}
\haidong{7} \end{ex}

%\newpage
\begin{ex}%book%[1P2Y1-1][Loai1]
Cho hai lực đồng quy có độ lớn $F_1=16\ N$ và $F_2=12\ N$. Hợp lực của chúng có thể có độ lớn $30\ N$ hoặc $3,5\ N$ được không?
\choice
{Có, nếu góc giữa hai vectơ lực có giá trị thích hợp}
{Có, nhưng chỉ trong một số trường hợp đặc biệt}
{\True Không, vì độ lớn hợp lực phải thỏa mãn $|F_1-F_2|\le F\le F_1+F_2$}
{Không, vì góc hợp bởi $2$ vectơ lực không thể bằng $0^\circ$ hoặc $180^\circ$}
\loigiai{
\begin{itemize}
\item Độ lớn hợp lực nhỏ nhất khi $\left({{F_1 } ,{F_2}}\right)=180^\circ$. 
\item Khi đó: $F_{\min}=| F_1-F_2 |=| 16- 12 |=4$ N
\item Độ lớn hợp lực lớn nhất khi $\left({{F_1 } ,{F_2}}\right)=0^\circ$.
\item Khi đó: $F_{\max}=F_1+F_2=16+12=28$ N.
\item Vậy F không thể bằng 30 N hoặc 3,5 N được.
\end{itemize}	
}
\haidong{7} \end{ex}
\subsubsection*{PHẦN 2. Thí sinh trả lời từ câu 1 đến câu 4. Trong mỗi ý a), b), c), d) ở mỗi câu, thí sinh chọn đúng hoặc sai.}
\setcounter{ex}{0}
\begin{ex}%book
Lực và năng lượng là hai đại lượng quan trọng nhất trong lĩnh vực Vật lí.
\choice
{Nếu không có lực tác dụng vào vật thì vật không thể chuyển động được}
{\True Lực là nguyên nhân gây ra biến dạng hoặc phá hủy một vật}
{Lực là nguyên nhân duy trì chuyển động của một vật}
{\True Lực là nguyên nhân làm biến đổi chuyển động của vật}
\loigiai{
\begin{itemchoice}
\itemch
\itemch
\itemch
\itemch
\end{itemchoice}
} \haidong{7}
\end{ex}
\begin{ex}%book
Các yếu tố của lực bao gồm điểm đặt, phương, chiều và độ lớn.
\choiceTF[t]
{\True Lực là đại lượng vectơ}
{\True Có thể tổng hợp các lực đồng quy theo quy tắc hình bình hành}
{Lực là yếu tố cần thiết để duy trì chuyển động của vật}
{\True Lực tác dụng lên vật gây ra gia tốc cho vật}
\loigiai{
\begin{itemchoice}
\itemch
\itemch
\itemch
\itemch
\end{itemchoice}
} \haidong{7}
 \end{ex}
\begin{ex}%book
Trong cơ học cổ điển, lực là đại lượng vectơ đặc trưng cho sự tương tác giữa các vật, và là nguyên nhân gây ra gia tốc cho vật. 
\choiceTF[t]
{Nếu không có lực tác dụng vào vật thì vật không chuyển động được}
{Nếu thôi không tác dụng lực vào vật thì vật đang chuyển động sẽ dừng lại}
{\True Lực gây ra gia tốc cho vật}
{\True Nếu có duy nhất một lực tác dụng lên vật thì vận tốc của vật bị thay đổi}
\loigiai{
\begin{itemchoice}
\itemch
\itemch
\itemch
\itemch
\end{itemchoice}
} \haidong{7}
 \end{ex}
\begin{ex}%book
Trong cơ học, khi nghiên cứu tương tác giữa các vật, người ta thường sử dụng phép phân tích lực để đơn giản hóa bài toán.
\choiceTF[t]
{Phép phân tích lực tuân theo quy tắc hình bình hành}
{Phép phân tích lực là phép thay thế một lực bằng một lực giống hệt nó}
{\True Phép phân tích lực là phép thay thế một lực bằng hai hay nhiều lực thành phần gây ra tác dụng giống nó}
{\True Phép phân tích lực là phép làm ngược lại với phép tổng hợp lực}
\loigiai{
\begin{itemchoice}
\itemch
\itemch
\itemch
\itemch
\end{itemchoice}
} \haidong{7}
 \end{ex}
\newpage 
\TieudeT{PHONG TỎA VẬT LÍ}
\subsubsection*{PHẦN I. Thí sinh trả lời câu hỏi từ câu 1 đến câu 18. Mỗi câu hỏi thí sinh chỉ chọn một phương án}
\setcounter{ex}{0}
\begin{ex}%book%[1P2Y1-3][Loai1]
Hai lực cân bằng là hai lực
\choice
{có cùng độ lớn, cùng giá và ngược chiều}
{\True có cùng độ lớn, cùng giá, ngược chiều và cùng tác dụng đồng thời lên một vật}
{có cùng độ lớn, ngược chiều và cùng tác dụng đồng thời lên một vật}
{có cùng độ lớn, cùng giá, ngược chiều và cùng tác dụng lên một vật trong những khoảng thời gian như nhau}
\loigiai{
Hai lực cân bằng là hai lực có cùng độ lớn, cùng giá, ngược chiều và cùng tác dụng đồng thời lên một vật.
}
%<MyLT>
\haidong{8} \end{ex}
\begin{ex}%book%[1P2Y1-3][Loai1]
Hai lực cân bằng \textbf{không} thể có
\choice
{cùng phương}
{\True cùng hướng}
{cùng độ lớn}
{cùng giá}
\loigiai{
Hai lực cân bằng là hai lực cùng phương (giá), ngược chiều (hướng), có độ lớn bằng nhau.
}
\haidong{8} \end{ex}
\begin{ex}%book
Cặp lực cân bằng \textbf{không} có đặc điểm?
\choice
{Cùng tác dụng vào một vật}
{Cùng phương, ngược chiều}
{Cùng độ lớn}
{\True Cùng phương, cùng chiều}
\loigiai{Chọn D

}\haidong{8} \end{ex}

\begin{ex}%book
Phát biểu nào sau đây \textbf{sai} khi nói về trọng lực?
\choice
{Trọng lực xác định bởi biểu thức $\vec{P}=m . \vec{g}$}
{Trọng lực là lực hút của Trái Đất tác dụng lên vật}
{Trọng lực tác dụng lên vật tỉ lệ thuận với khối lượng vật}
{\True Tại một nơi trên Trái Đất trọng lực tác dụng lên vật tỉ lệ thuận với gia tốc rơi tự do}
\loigiai{
Chọn D

}
\haidong{7} \end{ex}
\begin{ex}%book
Trọng tâm của vật là điểm đặt của
\choice
{\True trọng lực tác dụng vào vật}
{lực đàn hồi tác dụng vào vật}
{lực hướng tâm tác dụng vào vật}
{lực từ trường Trái Đất tác dụng vào vật}
\loigiai{
Chọn A
}
\haidong{7} \end{ex}
\begin{ex}%book
Gia tốc rơi tự do trên Trái Đất là $9,8 \ m / s^2$ còn trên sao Hỏa là $3,7 \ m / s^2$. Nếu một nhà du hành vũ trụ từ Trái Đất lên sao Hỏa sẽ có
\choice
{khối lượng và trọng lượng đều giảm đi}
{khối lượng và trọng lượng không đổi}
{\True khối lượng không đổi còn trọng lượng giảm đi}
{khối lượng giảm đi còn trọng lượng không đổi}
\loigiai{
Chọn C

}
\haidong{7} \end{ex}


\begin{ex}%book
Tình huống nào dưới đây \textbf{không} liên quan đến lực căng?
\begin{center}
\begin{tikzpicture}
\node[] at (0,0){\includegraphics[width=15cm]{images/tr69}};
\end{tikzpicture}
\end{center}
\choice
{a}
{b}
{\True c}
{d}
\haidong{7} \end{ex}
\begin{ex}%book%[1P2Y1-3][Loai1]
Muốn cho một chất điểm đứng cân bằng thì hợp lực của các lực tác dụng lên nó phải
\choice
{cân bằng với trọng lực của vật}
{\True bằng không}
{không đổi theo thời gian}
{khác không}
\loigiai{
Muốn cho một chất điểm đứng cân bằng thì hợp lực của các lực tác dụng lên nó phải bằng 0.
}
\haidong{8} \end{ex}
\begin{ex}%book
Điều kiện nào sau đây đúng khi nói về cân bằng của vật rắn dưới tác dụng của hai lực?
\choice
{Hai lực tác dụng phải song song, ngược chiều}
{Hai lực tác dụng phải trực đối}
{\True Hai lực tác dụng phải bằng nhau, ngược chiều}
{Hai lực tác dụng phải bằng nhau}

\loigiai{Chọn C

}\haidong{8} \end{ex}
\begin{ex}%book
Một chất điểm chịu tác dụng của ba lực cân bằng khi hợp lực của hai lực có
\choice
{cùng giá, cùng chiều, cùng độ lớn với lực thứ 3 }
{\True cùng giá, ngược chiều, cùng độ lớn với lực thứ 3}
{cùng phương, ngược chiều, cùng độ lớn với lực thứ 3}
{cùng phương, cùng chiều, cùng độ lớn với lực thứ 3}
\loigiai{Chọn B
}\haidong{8} \end{ex}
\begin{ex}%book
Trọng lực tác dụng lên vật có
\choice
{độ lớn luôn thay đổi}
{\True điểm đặt tại trọng tâm của vật, phương thẳng đứng, chiều từ trên xuống}
{điểm đặt tại trọng tâm của vật, phương thẳng đứng, chiều từ dưới lên}
{điểm đặt bất kì trên vật, phương thẳng đứng, chiều từ trên xuống}
\loigiai{
Chọn B
}
\haidong{7} \end{ex}
\begin{ex}%book%[1P2B5-3][Loai1]
Một vật có khối lượng m đặt ở nơi có gia tốc trọng trường g. Phát biểu nào sau đây là \textbf{sai}?
\choice
{Trọng lực có độ lớn được xác định bởi biểu thức $P = mg$}
{Điểm đặt của trọng lực là trọng tâm của vật}
{\True Trọng lực tỉ lệ nghịch với khối lượng của vật}
{Trọng lực là lực hút của Trái Đất tác dụng lên vật}
\loigiai{
Trọng lực $P = mg \Rightarrow$ trọng lực tỉ lệ thuận với khối lượng của vật.
}
\haidong{7} \end{ex}

\begin{ex}%book%SBT CTST 11.2
Tại cùng một địa điểm, hai vật có khối lượng $m_1<m_2$, trọng lực tác dụng lên hai vật lần lượt là $P_1$ và $P_2$ luôn thoả mãn điều kiện
\choice
{$P_1=P_2$}
{$\dfrac{P_1}{P_2}<\dfrac{m_1}{m_2}$}
{$P_1>P_2$}
{\True $\dfrac{P_1}{P_2}=\dfrac{m_1}{m_2}$}
\haidong{7} \end{ex}
\begin{ex}%book
Gia tốc rơi tự do trên Trái Đất là $9,8 \ m / s^2$ còn trên Mặt Trăng là $1,6 \ m / s^2$. Nếu một nhà du hành vũ trụ từ Mặt Trăng trở lại Trái Đất thì
\choice
{khối lượng và trọng lượng đều tăng lên}
{khối lượng và trọng lượng không đổi}
{\True khối lượng không đổi còn trọng lượng tăng xấp xỉ 6 lần}
{khối lượng không đổi còn trọng lượng giảm xấp xỉ 6 lần}
\loigiai{

Chọn C
$\dfrac{P_{T D}}{P_{M T}}=\dfrac{g_{T D}}{g_{M T}}=\dfrac{9,8}{1,6}=6,125$

}
\haidong{7} \end{ex}
\begin{ex}%book%SBT CTST 11.7
Khi vật treo trên sợi dây nhẹ cân bằng thì trọng lực tác dụng lên vật
\choice
{cùng hướng với lực căng dây}
{\True cân bằng với lực căng dây}
{hợp với lực căng dây một góc $90^\circ$}
{bằng không}
\haidong{7} \end{ex}
\begin{ex}%book
Đơn vị lực căng dây là
\choice
{Watts $(W)$}
{Joules (J)}
{\True Newton $(N)$}
{Radians (Rad)}
\loigiai{
Chọn C
}
\haidong{7} \end{ex}
\begin{ex}%Tailieuchuan
\immini[thm]{Dùng tay kéo một vật nặng như hình dưới dây. Lực căng dây tác dụng lên
\choice
{tay người}
{vật nặng}
{\True cả tay và vật nặng}
{điểm giữa của sợi dây}}{\includegraphics[width=4cm]{images/2025_04_24_746f9b43406b3514a111g-07(1)}}
\loigiai{}
\end{ex}
\begin{ex}%Tailieuchuan
\immini[thm]{Một chiếc nhẫn đồng chất tiết diện đều như hình vẽ. Trọng tâm của chiếc nhẫn nằm tại vị trí	
\choice
{$1$}
{\True $2$}
{$3$}
{$4$}}{\includegraphics[width=4cm]{images/2025_04_24_746f9b43406b3514a111g-08}
}	
\loigiai{}
\end{ex}























\subsubsection*{PHẦN 2. Thí sinh trả lời từ câu 1 đến câu 4. Trong mỗi ý a), b), c), d) ở mỗi câu, thí sinh chọn đúng hoặc sai.}
\setcounter{ex}{0}

\begin{ex}%book %[1P2B5-3][Loai1]
Khi một vật nằm gần bề mặt Trái Đất, lực hút của Trái Đất tác dụng lên vật được gọi là trọng lực.
\choiceTF[t]
{Trọng lực chính là tên gọi khác của trọng lượng}
{Trọng lực là lực hút của vật vào Trái Đất}
{\True Trọng lực là trường hợp riêng của lực hấp dẫn}
{Trọng lực là lực hút của Trái Đất vào các thiên thể}
\loigiai{ 
\begin{itemchoice}
\itemch 
\itemch 
\itemch 
\itemch 
\end{itemchoice}
}
\haidong{7} \end{ex}
\begin{ex}%book
Trọng lượng là lực do Trái Đất tác dụng lên vật và phụ thuộc vào gia tốc trọng trường tại nơi đặt vật.
\choiceTF[t]
{\True Trọng lượng là độ lớn trọng lực tác dụng lên vật}
{Trọng lượng của vật luôn không đổi tại mọi vị trí trên Trái Đất}
{\True Trọng lượng được tính theo công thức $P=mg$}
{\True Trọng lượng có thể được đo bằng lực kế}
\loigiai{
\begin{itemchoice}
\itemch 
\itemch 
\itemch 
\itemch 
\end{itemchoice}
}
\haidong{7} \end{ex}
\begin{ex}
Hai lực đồng quy $\vec{F_1}$, $\vec{F_2}$ có độ lớn lần lượt là $6\ N$ và $8\ N$.
\choiceTF[t]
{\True Độ lớn của hợp lực là $2\ N$, khi góc hợp bởi hai lực là $180^\circ$}
{Độ lớn của hợp lực là $4\ N$, khi góc hợp bởi hai lực là $60^\circ$}
{\True Độ lớn của hợp lực là $10\ N$, khi $\vec{F_1}$ và  $\vec{F_2}$ vuông góc}
{Hợp lực có giá trị lớn nhất bằng $16\ N$}
\loigiai{
\begin{itemchoice}
\itemch
\itemch
\itemch
\itemch
\end{itemchoice}}
\end{ex}
\begin{ex}
Cho hai lực đồng quy có độ lớn $F_1=16\ N$; $F_2=12\ N$.
\choiceTF[t]
{\True Độ lớn của hợp lực của hai lực này khi chúng hợp với nhau một góc $\alpha=0^\circ $ là $28\ N$}
{\True Hợp lực có độ lớn $20\ N$ thì góc giữa hai lực là $90^\circ$}
{Hai lực ngược chiều nhau, hợp lực có độ lớn $28\ N$}
{Hợp hai lực luôn lớn hơn $28\ N$}
\loigiai{
\begin{itemchoice}
\itemch
\itemch
\itemch
\itemch
\end{itemchoice}}
\end{ex}





















